%================================================
% HƯỚNG 2 — SO SÁNH TỈNH/VÙNG
%================================================
\section{So sánh khác biệt giữa tỉnh và vùng (Hướng 2)}

Hướng phân tích này trả lời ba câu hỏi: (i) vùng và tỉnh nào khác biệt trong
cùng một năm, (ii) có tỉnh nào là ngoại lệ thống kê trong phân bố điểm của năm
đó, và (iii) tỉnh nào thay đổi nhiều nhất giữa hai mốc do người dùng chọn.
Nguồn dữ liệu là Chỉ số Hiệu quả Quản trị và Hành chính công cấp tỉnh ở Việt
Nam (PAPI) do UNDP, CECODES và RTA công bố \cite{papi_vietnam}.

Dashboard chỉ đọc dữ liệu đã xử lý theo cấu trúc tỉnh--năm. Thước đo ``tổng 6
lĩnh vực gốc'' được dùng liên tục cho giai đoạn 2011--2024; thước đo tổng đủ 8
lĩnh vực chỉ dùng từ 2018. Hai thước đo không được so sánh chéo. Các quan sát
không có điểm cho thước đo đang chọn được ghi nhận là thiếu và không nội suy.

\subsection{Phương pháp và tương tác}

Ở một năm và một thước đo, dashboard vẽ choropleth 63 tỉnh, boxplot theo sáu
vùng, bảng xếp hạng top/bottom, và profile radar của tỉnh được chọn so với
trung bình vùng và toàn bộ tỉnh có dữ liệu. Người dùng có thể chọn tỉnh bằng
click trên bản đồ hoặc bằng selectbox; hai cách đều cập nhật cùng một benchmark.

Ngoại lệ được tính bằng z-score mẫu trong đúng snapshot năm đang xem,
\(z_i=(x_i-\bar{x})/s\), và chỉ gắn cờ khi \(|z_i|>2\). Nếu ít hơn hai quan
sát hoặc độ lệch chuẩn bằng 0, hệ thống không suy diễn ngoại lệ. Slopegraph chỉ
nối tỉnh có đủ điểm ở cả hai năm được chọn; vì vậy một đường luôn thể hiện một
so sánh quan sát được, không phải giá trị bù thiếu.

\subsection{Ví dụ số liệu tái tạo}

Với tổng 6 lĩnh vực gốc năm 2024, dữ liệu có 61 tỉnh đủ điểm. Trung bình vùng
cao nhất là Đồng bằng sông Hồng (37.29) và thấp nhất là Tây Nguyên (34.69),
chênh 2.60 điểm. Ba tỉnh điểm cao nhất là Quảng Ninh (39.85), Bình Thuận
(39.65) và Tây Ninh (39.54); ba điểm thấp nhất là Cần Thơ (33.19), Kiên Giang
(33.47) và Đắk Nông (33.96). Các số này được tái tạo trực tiếp từ
\texttt{data/processed/agg\_province\_year.parquet} bằng helper
\texttt{analysis.provincial}, không chép từ nhánh phát triển cũ.

\subsection{Giới hạn}

Kết quả mô tả chênh lệch điểm khảo sát, không suy ra quan hệ nhân quả hay xếp
hạng hành chính ngoài phạm vi dữ liệu hiện có. Số tỉnh trong từng biểu đồ có thể
thấp hơn 63 khi một tỉnh thiếu điểm ở năm/thước đo được chọn; dashboard luôn
hiển thị phạm vi này cùng nguồn PAPI.
